\subsection{Symmetric hilbertian powers}

Let E be a hilbertian space, and let n be a positive integer. Let \(\hat{\mathbf{T}}^{n}(\mathrm{E})\) or \(E^{\otimes_{n}}\) denote the tensor product of n hilbertian spaces each equal to E. In other words, \(\hat{\mathbf{T}}^{n}(\mathrm{E})\) is the completion of the space \(\mathbf{T}^{n}(\mathrm{E}) = \mathrm{E} \otimes \ldots \otimes \mathrm{E}\) (n factors) for the Hausdorff prehilbertian space structure defined by

\[
\left\langle x _ {1} \otimes \dots \otimes x _ {n} \mid y _ {1} \otimes \dots \otimes y _ {n} \right\rangle = \prod_ {i = 1} ^ {n} \left\langle x _ {i} \mid y _ {i} \right\rangle .
\]

If \((e_i)_{i\in I}\) is an orthonormal basis of \(E\) , the family of vectors \(e_{i_1} \otimes \ldots \otimes e_{i_n}\) for \(i_1, \ldots, i_n\) in \(I\) , is an orthonormal basis of \(\hat{\mathbf{T}}^n(E)\) (V, p.~29, cor. 1). We get \(\hat{\mathbf{T}}^\circ(E) = K\) .

Let \(\sigma \in \mathfrak{S}_n\) be a permutation of the set \(\{1, 2, ..., n\}\) . By V, p.~27, there exists an automorphism \(p_{\sigma}\) of \(\hat{\mathbf{T}}^n(\mathbf{E})\) characterized by

\[
p _ {\sigma} (x _ {1} \otimes \dots \otimes x _ {n}) = x _ {\sigma^ {- 1} (1)} \otimes \dots \otimes x _ {\sigma^ {- 1} (n)}.
\]

We have \(p_{\sigma \tau} = p_{\sigma}p_{\tau}\) for \(\sigma, \tau\) in \(\mathfrak{S}_n\) , and consequently the endomorphism \(\Pi_n = \frac{1}{n!}\sum_{\sigma \in \mathfrak{S}_n}p_\sigma\) of the vector space \(\hat{\mathbf{T}}^n (\mathrm{E})\) is the orthoprojector onto the subspace of all elements left invariant by \(\mathfrak{S}_n\) . However (A, III, § 6, No.~3), \(\Pi_n\) maps the « algebraic » tensor product \(\mathbf{T}^n (\mathrm{E})\) onto the subspace \(\mathbf{T}\mathbf{S}^n (\mathrm{E})\) of all symmetric tensors of order \(n\) . In other words, the image of \(\Pi_n\) is the completion of the space \(\mathbf{T}\mathbf{S}^n (\mathrm{E})\) endowed with a scalar product induced by that of \(\mathbf{T}^n (\mathrm{E})\) ; this completion will be denoted by \(\widehat{\mathbf{T}\mathbf{S}}^n (\mathrm{E})\) .

Let \(\mathbf{S}^n (\mathrm{E})\) be the \(n\) th symmetric power of the vector space E (A, III, § 6, No.~1). The canonical mapping from \(\mathbf{T}^n (\mathrm{E})\) onto \(\mathbf{S}^n (\mathrm{E})\) defines by restriction an isomorphism \(\lambda_{n}\) from \(\mathbf{TS}^n (\mathrm{E})\) onto \(\mathbf{S}^n (\mathrm{E})\) . We verify immediately that the inverse isomorphism is given by

\[
\mu_ {n} (x _ {1} \dots x _ {n}) = \Pi_ {n} (x _ {1} \otimes \dots \otimes x _ {n}) = \frac {1}{n !} \sum_ {\sigma \in \mathfrak {S} _ {n}} x _ {\sigma^ {- 1} (1)} \otimes \dots \otimes x _ {\sigma^ {- 1} (n)}\tag{13}
\]

for \(x_{1},\ldots ,x_{n}\) in E.

We define a Hausdorff prehilbertian space structure on \(\mathbf{S}^{n}(\mathrm{E})\) by putting

\[
\langle u | v \rangle = n! \left\langle \mu_ {n} (u) | \mu_ {n} (v) \right\rangle .\tag{14}
\]

We then have (compare with formula (29) of A, III, § 11, No.~5)

\[
\left\langle x _ {1} \dots x _ {n} \mid y _ {1} \dots y _ {n} \right\rangle = \sum_ {\sigma \in \mathfrak {S} _ {n}} \prod_ {i = 1} ^ {n} \left\langle x _ {i} \mid y _ {\sigma (i)} \right\rangle ,\tag{15}
\]

and in particular

\[
\langle x ^ {n} | y ^ {n} \rangle = n! \langle x | y \rangle^ {n}.\tag{16}
\]

Let \(\hat{\mathbf{S}}^n (\mathrm{E})\) denote the completion of the pre-Hilbertian space \(\mathbf{S}^n (\mathrm{E})\) and \(\hat{\mathbf{S}} (\mathrm{E})\) the external hilbertian sum of the hilbertian spaces \(\hat{\mathbf{S}}^n (\mathrm{E})\) . We can show (V, p.~73, exerc. 1) that the multiplication in the algebra \(\mathbf{S}(\mathrm{E})\) cannot be extended by continuity to \(\mathbf{S}(\mathrm{E})\) , unless E is just 0.

PROPOSITION 4. --- Let \((e_{i})_{i\in I}\) be an orthonormal basis of the hilbertian space E. For every \(\alpha\) in \(\mathbf{N}^{(1)}\) , put

\[
z _ {\alpha} = \prod_ {i \in I} e _ {i} ^ {\alpha_ {i}} / (\alpha_ {i}!) ^ {1 / 2}.\tag{17}
\]

Then \((z_{\alpha})_{\alpha \in \mathbf{N}^{(1)}}\) is an orthonormal basis of \(\hat{\mathbf{S}} (\mathrm{E})\) .

Let \(E_{0}\) be the vector subspace of E generated by the vectors \(e_{i}\) for i ranging over I. Then the \(z_{\alpha}\) form a basis of the vector space \(\mathbf{S}(E_{0})\) (A, III, § 6, No.~6). But \(E_{0}\) is dense in E, and the multilinear mapping \((x_{1},...,x_{n})\mapsto x_{1}\ldots x_{n}\) from \(E\times\cdots\times E\) into \(\mathbf{S}(E)\) is continuous for all \(n\geqslant1\) ; hence \(\mathbf{S}(E_{0})\) is dense in \(\mathbf{S}(E)\) . It is now enough to prove that the family of the \(z_{\alpha}\) is orthonormal. First observe that \(\hat{\mathbf{S}}^{n}(\mathrm{E})\) and \(\hat{\mathbf{S}}^{m}(\mathrm{E})\) are orthogonal for \(n \neq m\) . Hence it is enough to prove the formula

\[
\langle z _ {\alpha} | z _ {\beta} \rangle = \left\{ \begin{array}{l l} 1 & \text {if} \quad \alpha = \beta \\ 0 & \text {if} \quad \alpha \neq \beta \end{array} \right.
\]

when \(|\alpha| = \sum_{i \in I} \alpha_i\) and \(|\beta| = \sum_{i \in I} \beta_i\) are equal to the same integer \(n\) .

Consider a partition \((\mathbf{P}_i)_{i\in I}\) of the set \(\{1,2,\dots,n\}\) such that \(\operatorname{Card} \mathbf{P}_i = \alpha_i\) for all \(i\in I\) . Put \(x_k = e_i\) if \(k\) belongs to \(\mathbf{P}_i\) , then \(x_1\ldots x_n = \prod_{i\in I}e_i^{\alpha_i}\) . Similarly we define \((\mathbf{Q}_i)_{i\in I}\) and \(y_k\) in such a way that \(\operatorname{Card} \mathbf{Q}_i = \beta_i\) and \(y_1\ldots y_n = \prod_{i\in I}e_i^{\beta_i}\) . Since the \(e_i\) are mutually orthogonal, we have \(\langle x_k|y_{\sigma(k)}\rangle = 0\) except if there exists an indice \(i\in I\) such that \(k\in \mathbf{P}_i\) and \(\sigma (k)\in \mathbf{Q}_i\) . By formula (15), we then have \(\langle x_1\ldots x_n|y_1\ldots y_n\rangle = 0\) unless there exists a permutation \(\sigma \in \mathfrak{S}_n\) such that \(\sigma (\mathbf{P}_i) = \mathbf{Q}_i\) for all \(i\in I\) , which implies that \(\alpha = \beta\) . Then \(\langle z_{\alpha}|z_{\beta}\rangle = 0\) for \(\alpha \neq \beta\) . The same argument proves that \(\| x_1\ldots x_n\|^2\) is equal to the number of the \(\sigma \in \mathfrak{S}_n\) such that \(\sigma (\mathbf{P}_i) = \mathbf{P}_i\) for all \(i\in I\) , hence equal to \(\prod_{i=1}^{\infty}\alpha_i!\) . We get \(\| z_\alpha \| = 1\) , and the proposition is proved.

COROLLARY. --- Suppose that the hilbertian space E is the direct sum of the orthogonal subspaces M and N. The canonical isomorphism g from \(\mathbf{S}(\mathbf{M}) \otimes \mathbf{S}(\mathbf{N})\) onto \(\mathbf{S}(\dot{\mathbf{E}})\) (A, III, §6, No.~6) extends uniquely to a hilbertian space isomorphism h from \(\hat{\mathbf{S}}(\mathbf{M}) \otimes_{2} \hat{\mathbf{S}}(\mathbf{N})\) onto \(\hat{\mathbf{S}}(\mathbf{E})\) .

Let \((e_{i})_{i\in I}\) (resp. \((f_{j})_{j\in J}\) ) be an orthonormal basis of the hilbertian space M (resp. N) and let \(M_{0}\) (resp. \(N_{0}\) ) be the vector subspace of E generated by the vectors \(e_{i}\) (resp. \(f_{j}\) ). Put \(E_{0}=M_{0}+N_{0}\) and let \(g_{0}\) be the canonical isomorphism from \(S(M_{0})\otimes S(N_{0})\) onto \(S(E_{0})\) . Put

\[
z _ {\alpha} = \prod_ {i \in \mathbf {I}} e _ {i} ^ {\alpha_ {i}} / (\alpha_ {i}!) ^ {1 / 2}, t _ {\beta} = \prod_ {j \in \mathbf {J}} f _ {j} ^ {\beta _ {j}} / (\beta_ {j}!) ^ {1 / 2}
\]

for \(\alpha \in \mathbf{N}^{(I)}\) and \(\beta \in \mathbf{N}^{(J)}\) . By prop. 4, we have thus defined the orthonormal bases \((z_{\alpha})_{\alpha \in \mathbf{N}^{(I)}}\) for \(\hat{\mathbf{S}}(\mathbf{M})\) , \((t_{\beta})_{\beta \in \mathbf{N}^{(J)}}\) for \(\hat{\mathbf{S}}(\mathbf{N})\) and \((z_{\alpha}t_{\beta})_{\alpha \in \mathbf{N}^{(I)},\beta \in \mathbf{N}^{(J)}}\) for \(\hat{\mathbf{S}}(\mathbf{E})\) . Since we have \(z_{\alpha}t_{\beta} = g_0(z_{\alpha}\otimes t_{\beta})\) , and since the elements \(z_{\alpha}\otimes t_{\beta}\) form an orthonormal basis of \(\hat{\mathbf{S}}(\mathbf{M})\hat{\otimes}_2\hat{\mathbf{S}}(\mathbf{N})\) (V, p.~29, cor. 1), we see that \(g_0\) extends to a hilbertian space isomorphism \(h:\hat{\mathbf{S}}(\mathbf{M})\hat{\otimes}_2\hat{\mathbf{S}}(\mathbf{N})\to \hat{\mathbf{S}}(\mathbf{E})\) . By the construction, we have

\[
h (x _ {1} \dots x _ {m} \otimes y _ {1} \dots y _ {n}) = x _ {1} \dots x _ {m} y _ {1} \dots y _ {n}
\]

for all vectors \(x_{1}, \ldots, x_{m}\) in \(M_{0}\) and \(y_{1}, \ldots, y_{n}\) in \(N_{0}\) . By continuity, the same relation also holds for the vectors \(x_{1}, \ldots, x_{m}\) in M and the vectors \(y_{1}, \ldots, y_{n}\) in N; in other words, h extends g. The uniqueness of h is clear.

Let E and F be two hilbertian spaces and \(u \in \mathcal{L}(E; F)\) . The linear mapping \(\hat{\mathbf{T}}^{n}(u) = u \otimes_{2} \ldots \otimes_{2} u\) (n factors) from \(\hat{\mathbf{T}}^{n}(E)\) into \(\hat{\mathbf{T}}^{n}(F)\) is continuous with norm \(\|u\|^{n}\) (V, p.~28, formula (10)). Moreover, formulas (13) and (14) of V, p.~30, show that there exists an isomorphism \(\phi_{n,\mathrm{E}}\) from \(\hat{\mathbf{S}}^n (\mathrm{E})\) onto the subspace \(\widehat{\mathbf{T}\mathbf{S}}^n (\mathrm{E})\) of \(\hat{\mathbf{T}}^n (\mathrm{E})\) , and only one, such that

\[
\phi_ {n, \mathrm{E}} (x _ {1} \dots x _ {n}) = \frac {1}{(n !) ^ {1 / 2}} \sum_ {\sigma \in \mathfrak {S} _ {n}} x _ {\sigma (1)} \otimes \dots \otimes x _ {\sigma (n)} \quad (x _ {1},..., x _ {n} \text { in } \mathrm{E}).\tag{18}
\]

Hence there exists a continuous linear mapping \(\hat{\mathbf{S}}^{n}(u)\) from \(\hat{\mathbf{S}}^{n}(\mathrm{E})\) into \(\hat{\mathbf{S}}^{n}(\mathrm{F})\) and only one which makes the following diagram commutative

\[
\begin{array}{c c c} \hat {\mathbf {S}} ^ {n} (\mathrm{E}) & \xrightarrow {\phi_ {n , \mathrm{E}}} & \hat {\mathbf {T}} ^ {n} (\mathrm{E}) \\ \hat {\mathbf {S}} ^ {n} (u) \Big \downarrow & & \Big \downarrow \hat {\mathbf {T}} ^ {n} (u) \\ \hat {\mathbf {S}} ^ {n} (\mathrm{F}) & \xrightarrow {\phi_ {n , \mathrm{F}}} & \hat {\mathbf {T}} ^ {n} (\mathrm{F}) \end{array}
\]

We now prove the formula

\[
\left\| \hat {\mathbf {S}} ^ {n} (u) \right\| = \| u \| ^ {n}.\tag{19}
\]

We clearly have \(\| \hat{\mathbf{S}}^n (u)\| \leqslant \| \hat{\mathbf{T}}^n (u)\| = \| u\| ^n\) . Further, for all \(x\in \mathrm{E}\) , we have \(\hat{\mathbf{S}}^n (u)(x^n) = u(x)^n,\| x^n\| = (n!)^{1 / 2}\| x\| ^n\) and \(\| u(x)^n\| = (n!)^{1 / 2}\| u(x)\| ^n\) , which gives

\[
\left\| \hat {\mathbf {S}} ^ {n} (u) \right\| \| x \| ^ {n} \geqslant \left\| u (x) \right\| ^ {n};
\]

it follows immediately that \(\| \hat{\mathbf{S}}^n (u)\| \geqslant \| u\| ^n\) , hence formula (19).

It is clear that we have the formulas

\[
\hat {\mathbf {S}} ^ {n} (1 _ {\mathrm{E}}) = 1 _ {\hat {\mathbf {S}} ^ {n} (\mathrm{E})}\tag{20}
\]

\[
\hat {\mathbf {S}} ^ {n} (v \circ u) = \hat {\mathbf {S}} ^ {n} (v) \circ \hat {\mathbf {S}} ^ {n} (u) \quad \text { for } \quad v \in \mathscr {L} (\mathrm{F}; \mathrm{G}).\tag{21}
\]

Finally, \(\hat{\mathbf{S}}^n (u)\) coincides on \(\mathbf{S}^n (\mathrm{E})\) with the linear mapping \(\mathbf{S}^n (u):\mathbf{S}^n (\mathrm{E})\to \mathbf{S}^n (\mathrm{F})\) defined in A, III, § 6, No.~2 since it transforms \(x_{1}\ldots x_{n}\) into \(u(x_{1})\ldots u(x_{n})\) for every \(x_{1},\ldots ,x_{n}\) in E.

Examples. --- * 1) Let \(d \geqslant 1\) be an integer and \(\omega\) a positive function on \(\mathbf{R}^{d}\) , locally integrable with respect to the Lebesgue measure \(\mu\) . Let E be the hilbertian space \(\mathrm{L}^{2}(\mathbf{R}^{d}, \omega, \mu)\) , and let \(\mathbf{S} = \mathbf{S}(\mathrm{E})\) . Then \(\mathbf{S}\) can be identified with the space of all sequences \(f = (f_{n})_{n \geqslant 0}\) , where each \(f_{n}\) is a function on \((\mathbf{R}^{d})^{n}\) which is measurable with respect to the Lebesgue measure \(\mu \otimes \ldots \otimes \mu\) ( \(n\) factors) and invariant under the permutations of the \(n\) factors in \((\mathbf{R}^{d})^{n}\) , and is such that

\[
\| \mathbf {f} \| ^ {2} = \sum_ {n = 0} ^ {\infty} n! \int_ {\mathbf {R} ^ {d}} \dots \int_ {\mathbf {R} ^ {d}} \left| f _ {n} (\mathbf {x} _ {1}, \dots , \mathbf {x} _ {n}) \right| ^ {2} \omega (\mathbf {x} _ {1}) \dots \omega (\mathbf {x} _ {n}) d \mathbf {x} _ {1} \dots d \mathbf {x} _ {n}\tag{22}
\]

is finite. The norm \(\| \mathbf{f}\|\) in \(\mathbf{S}\) is defined by formula (22). The hilbertian space \(\mathbf{S}\) defined above is called the symmetric Fock space corresponding to the weight \(\omega_{*}\)

* 2) Let X be a Hausdorff topological space, \(\mu\) a positive measure of norm 1 on X and E a hilbertian subspace of the real hilbertian space \(L_{\mathbb{R}}^{2}(X,\mu)\) . We say that E is a gaussian space if the following equivalent conditions are satisfied:

\begin{enumerate}
\def\labelenumi{\alph{enumi})}
\item
  for all \(f \in \mathrm{E}\) , we have \(\int_{\mathbf{x}} e^{if} d\mu = \exp(-\|f\|^2/2)\) ;
\item
  for all \(f \in E\) with norm 1, the image of the measure \(\mu\) under f is the measure
\end{enumerate}

\[
(2 \pi) ^ {- 1 / 2} e ^ {- x ^ {2} / 2} d x.
\]

on \(\mathbf{R}\) .

Suppose E is a gaussian space. Let \(f_1, \ldots, f_n\) be functions whose classes \(f_i\) belong to E. We define a function: \(f_1 \ldots f_n\) : on X (called « Wick's product » of \(f_1, \ldots, f_n\) ) by the formula

\[
: f _ {1} \dots f _ {n} := \sum_ {0 \leqslant 2 p \leqslant n} (- 1) ^ {p} \sum_ {\sigma \in I _ {p}} \prod_ {i = 1} ^ {p} \left\langle f _ {\sigma (2 i - 1)} \mid f _ {\sigma (2 i)} \right\rangle \prod_ {j = 2 p + 1} ^ {n} f _ {\sigma (j)},\tag{23}
\]

where \(\mathbf{I}_p\) is the set of permutations \(\sigma\) of \(\{1,2,\dots,n\}\) such that we have

\[
\begin{array}{l} \sigma (1) <   \sigma (2), \dots , \sigma (2 p - 1) <   \sigma (2 p) \\ \sigma (1) <   \sigma (3) <   \dots <   \sigma (2 p - 1) \\ \sigma (2 p + 1) <   \sigma (2 p + 2) <   \dots <   \sigma (n). \end{array}
\]

Then there exists an isomorphism \(\phi\) from \(\hat{\mathbf{S}}(\mathrm{E})\) onto a hilbertian subspace of \(\mathrm{L}_{\mathbb{R}}^{2}(\mathrm{X},\mu)\) which transforms the product \(\dot{f}_1\dots \dot{f}_n\) of \(\dot{f}_1,\dots ,\dot{f}_n\) , calculated in \(\hat{\mathbf{S}}(\mathrm{E})\) , into \((:f_1\dots f_n:)\). Suppose that X is a Souslin space and that there exists a countable family \((f_{n})\) of functions whose classes belong to E and which separate the points of X. Then \(\phi\) is an isomorphism from \(\hat{\mathbf{S}}(\mathrm{E})\) onto \(\mathrm{L}_{\mathbb{R}}^{2}(\mathrm{X},\mu)\) .

